\documentclass[aspectratio=169]{beamer}
\usetheme{Madrid}
\usecolortheme{default}
\usepackage{amsmath}
\usepackage{graphicx}
\usepackage{tikz-cd}
\usepackage{multicol}

\setlength{\parindent}{0pt}
\setlength{\parskip}{1\baselineskip}

\title[Lecture 03]{An Introduction to E-Commerce\\\small with a Focus on Machine Learning, Deep Learning, and Artificial Intelligence}
\subtitle{Lecture 03: E-Commerce Data Foundations: Events, Experimentation, and Metrics}
\author{Dr. Haitham A. El-Ghareeb}
\institute{Faculty of Computers and Information Sciences \\ Mansoura University \\ Egypt}
\date{\today}

\begin{document}

\begin{frame}
  \titlepage
\end{frame}

\begin{frame}{Session plan (90 minutes)}
  \begin{enumerate}
    \item Why data foundations matter
    \item From business process to event model
    \item Event schema design
    \item Identity resolution and entity modeling
    \item Data quality as a product requirement
    \item Hands-on lab: build a minimal metrics layer
  \end{enumerate}
\end{frame}

\begin{frame}{Learning objectives}
  By the end of this lecture, you should be able to:
  \begin{itemize}
    \item Design an event-driven data model for storefront, marketplace, and mobile commerce scenarios.
    \item Distinguish identifiers, identities, sessions, and entities across operational and analytical systems.
    \item Explain how attribution, data quality, and metric definitions affect decision-making and ML performance.
    \item Plan controlled experiments and understand when adaptive approaches such as bandits are appropriate.
  \end{itemize}
\end{frame}

\section{Why data foundations matter}

\begin{frame}{Why data foundations matter}
  \begin{itemize}
    \item Product analytics and funnel analysis
    \item Operational reporting across commerce and ERP systems
    \item Reliable feature generation for ML models
    \item Controlled online experimentation
    \item Governance, auditing, and cross-functional alignment
  \end{itemize}
\end{frame}

\section{From business process to event model}

\begin{frame}{From business process to event model}
  \begin{itemize}
    \item \textbf{Events versus state:} \textbf{State} answers "what is true now?" Examples: current inventory, current product price, current loyalty tier.; \textbf{Events} answer "what happened?" Examples: product viewed, cart updated, payment failed, order shipped.
    \item \textbf{Event taxonomies:} \textbf{Discovery events:} homepage viewed, category viewed, search submitted, filter applied, result clicked.; \textbf{Consideration events:} product viewed, review expanded, comparison opened, wishlist updated.
  \end{itemize}
\end{frame}

\section{Event schema design}

\begin{frame}{Event schema design}
  \begin{itemize}
    \item \textbf{Required design principles:} \textbf{Business semantics before UI semantics:} schemas should reflect domain meaning, not implementation accidents.; \textbf{Stable field names:} renaming fields casually breaks dashboards, pipelines, and models.
    \item \textbf{Example event record:} The following conceptual record shows how a product view may be represented: The example illustrates an important practice: authenticated and anonymous identifiers may coexist.
  \end{itemize}
\end{frame}

\section{Identity resolution and entity modeling}

\begin{frame}{Identity resolution and entity modeling}
  \begin{itemize}
    \item \textbf{Identifiers are not interchangeable:} \textbf{User/account ID:} application-level identity for authentication and personalization.; \textbf{Customer/partner ID:} commercial identity used in ERP, billing, and support.
    \item \textbf{Identity stitching:} linking a guest cart to a logged-in account after authentication,; attaching a mobile device history to a user after sign-in,
    \item \textbf{Privacy-aware identity design:} minimizing collection of directly identifying information,; separating analytical IDs from operational personal data where possible,
  \end{itemize}
\end{frame}

\section{Data quality as a product requirement}

\begin{frame}{Data quality as a product requirement}
  \begin{itemize}
    \item \textbf{Data contracts:} event name and schema,; field definitions and allowed values,
    \item \textbf{Validation and observability:} row-count anomalies by event family,; null-rate spikes in required fields,
  \end{itemize}
\end{frame}

\section{Attribution and its limits}

\begin{frame}{Attribution and its limits}
  \begin{itemize}
    \item \textbf{Common attribution models:} \textbf{Last-touch attribution:} credits the final recorded touchpoint before conversion.; \textbf{First-touch attribution:} credits the earliest touchpoint in the conversion path.
    \item \textbf{Selection bias and measurement bias:} \textbf{Selection bias:} users exposed to an experience may differ systematically from users who are not.; \textbf{Survivorship bias:} only recorded users or channels appear in the analysis.
  \end{itemize}
\end{frame}

\section{Metrics layers and decision quality}

\begin{frame}{Metrics layers and decision quality}
  \begin{itemize}
    \item \textbf{North Star and supporting metrics:} traffic volume,; product-detail engagement,
    \item \textbf{Metric dimensions:} channel (web, app, store-assisted, marketplace),; customer segment (new, returning, VIP, B2B account tier),
    \item \textbf{Retention and cohort logic:} Did the user return to browse?; Did the user make another purchase?
  \end{itemize}
\end{frame}

\section{Experimentation in e-commerce}

\begin{frame}{Experimentation in e-commerce}
  \begin{itemize}
    \item \textbf{A/B testing fundamentals:} \textbf{Unit of randomization:} user, session, device, seller, or region.; \textbf{Exposure logic:} where and when eligibility is checked.
    \item \textbf{Common experimental pitfalls:} \textbf{Contamination:} the same user sees both variants through multiple devices or sessions.; \textbf{Novelty effects:} short-term excitement masks weak long-term performance.
    \item \textbf{Bandits and adaptive experimentation:} promotions on fast-moving catalogs,; ranking widgets with many interchangeable treatments,
  \end{itemize}
\end{frame}

\section{From analytics to machine learning features}

\begin{frame}{From analytics to machine learning features}
  \begin{itemize}
    \item recent category views for recommendation,
    \item failed login or payment attempts for fraud detection,
    \item coupon usage history for promotion targeting,
    \item time since last order for churn models,
    \item return behavior and order defects for seller-quality models.
  \end{itemize}
\end{frame}

\section{Hands-on lab: build a minimal metrics layer}

\begin{frame}{Hands-on lab: build a minimal metrics layer}
  \begin{itemize}
    \item \textbf{Lab scenario:} Assume a storefront emits five core events: , , , , and .
    \item \textbf{Lab tasks:} Create a canonical event table with required fields, event validation flags, and ingestion metadata.; Build daily metrics for sessions, product views, add-to-cart rate, checkout-start rate, purchase conversion rate, and day-7 repeat purchase.
    \item \textbf{Expected learning outcome:} By the end of the lab, students should be able to explain why a metric is not just an aggregation, but the result of event design choices, identity assumptions, and validation rules.
  \end{itemize}
\end{frame}

\section{Managerial and architectural implications}

\begin{frame}{Managerial and architectural implications}
  \begin{itemize}
    \item Which identifiers are authoritative, and where are they mastered?
    \item Which events belong in the clickstream versus operational event streams?
    \item How are storefront events linked to ERP objects such as orders, invoices, and shipments?
    \item What happens when a schema changes in the storefront while ERP integrations remain stable, or vice versa?
    \item Who approves metric definitions that drive executive dashboards and model targets?
  \end{itemize}
\end{frame}

\end{document}
